\documentclass[11pt,a4paper]{article}
\usepackage[margin=25mm]{geometry}
\usepackage{booktabs}
\usepackage{enumitem}
\usepackage{fancyhdr}
\usepackage[hidelinks]{hyperref}

\pagestyle{fancy}
\fancyhf{}
\fancyhead[L]{\small Oriole Packaging Systems}
\fancyhead[R]{\small ORI-MAN-PX400-Rev7}
\fancyfoot[C]{\small \thepage}

\title{PX-400 Field Service Manual\\ \large Fault E-217 Troubleshooting and Diagnostic D-17}
\author{Oriole Packaging Systems --- Service Engineering}
\date{Revision 7 \quad September 2026}

\begin{document}
\maketitle

\section{Safety}
Isolate mains power and lock out the line before touching the encoder connector,
harness, or controller cabinet. Verify zero energy with a meter. Never reseat
connectors with power applied. Record every intervention in the service work
order before leaving the site.

\section{Fault code reference}
\begin{table}[h]
\centering
\begin{tabular}{lll}
\toprule
Code & Meaning & First action \\
\midrule
E-210 & Index overtravel & Check mechanical stops, then home offset \\
E-217 & Encoder signal loss during index & See Section~3, do not replace parts first \\
E-223 & Controller watchdog timeout & Capture log, cycle control power \\
E-301 & Firmware image mismatch & Verify firmware version, reflash per release note \\
\bottomrule
\end{tabular}
\end{table}

\section{E-217 troubleshooting procedure}
E-217 means the motion controller lost a valid encoder signal during an indexed
movement. After firmware~4.8 the diagnostic sampling window is tighter, so
marginal encoder connections that used to recover now raise a stop. Work the
steps in order and stop at the first step that resolves the fault.

\begin{enumerate}[leftmargin=*]
\item Read the controller fault log and note the E-217 timestamps. Compare them
with recent firmware or maintenance changes.
\item Isolate power. Inspect the encoder connector (J7), harness routing, shield
termination, and visible cable damage along the full cable run.
\item Reseat the J7 connector. Check the locking tab and confirm the shield
drain wire is terminated at the controller end only.
\item Re-apply power and run Diagnostic D-17 (Section~4). Record pass or fail
with the fault log attached to the work order.
\item If D-17 fails, replace the indicated component per Section~5 and rerun D-17.
\item If D-17 passes repeatedly but E-217 stops continue, capture the controller
log and escalate to engineering with timestamps and the D-17 results.
\end{enumerate}

Do not downgrade firmware as the first response to E-217.

\section{Diagnostic D-17 procedure}
D-17 exercises the index axis through ten slow index cycles while sampling the
encoder channel each millisecond. A pass requires ten clean cycles with no
signal dropout longer than 2~ms. A fail reports the cycle number and the dropout
duration; attach this output to the work order as the diagnostic basis for any
part replacement.

\section{Replacement parts}
\begin{table}[h]
\centering
\begin{tabular}{lll}
\toprule
SKU & Description & Replace when \\
\midrule
EA-400 & PX-400 encoder assembly & D-17 fails after connector/harness check \\
HN-44 & Shielded encoder harness & Visible damage or shield fault found \\
PS-12 & Controller power module & Only on power-rail fault, never for E-217 alone \\
\bottomrule
\end{tabular}
\end{table}

\section{Connector and torque reference}
\begin{table}[h]
\centering
\begin{tabular}{ll}
\toprule
Item & Specification \\
\midrule
Encoder connector & J7, 17-pin circular, locking tab \\
J7 coupling torque & 1.2~N\,m \\
Cabinet ground lug & 2.5~N\,m \\
Shield termination & Controller end only, drain wire under clamp \\
\bottomrule
\end{tabular}
\end{table}

\end{document}
